% Three additional Beamer slides: GPU + parallel R
% Insert after the Rcpp/RcppParallel material or before the course summary.
% No images or additional LaTeX packages are required.

\subsection{GPU computing with R}

\begin{frame}{CPU parallelism and GPU acceleration}
\small
\begin{itemize}
  \item A CPU has a relatively small number of powerful, general-purpose cores; a GPU provides many simpler execution units for highly parallel numerical workloads.
  \item GPU acceleration is most useful for large matrix operations, tensor computations, deep-learning models, and large batches of similar operations.
  \item CPU parallelism is usually more suitable for irregular tasks, branching algorithms, many small independent jobs, and workflows dominated by disk or network access.
  \item A GPU does not accelerate ordinary R code automatically: computations must use a GPU-enabled library and data must be transferred to GPU memory.
\end{itemize}
\vspace{0.15cm}
\alert{Practical rule:} first create and verify a CPU baseline; use a GPU only when the workload is sufficiently large to compensate for setup and data-transfer costs.
\end{frame}

\begin{frame}[fragile]{GPU tensors in R with \texttt{torch}}
\begin{block}{The same tensor operation on the available device}
\begin{lstlisting}[language=R,basicstyle=\ttfamily\scriptsize]
library(torch)

use_gpu <- cuda_is_available()
device <- if (use_gpu) torch_device("cuda") else torch_device("cpu")
cat("Using:", device$type, "\n")

n <- 4000
x <- torch_randn(n, n, device = device)
y <- torch_randn(n, n, device = device)

z <- torch_matmul(x, y)
if (use_gpu) cuda_synchronize()

result <- z$mean()$item()
rm(x, y, z)
\end{lstlisting}
\end{block}
\begin{itemize}
  \item \texttt{cuda\_is\_available()} checks whether CUDA can be used.
  \item Keep intermediate tensors on the GPU; copy back only the final result needed in R.
  \item Synchronise before timing because GPU operations may execute asynchronously.
\end{itemize}
\end{frame}

\begin{frame}[fragile]{Requesting a GPU with Slurm}
\begin{block}{Example batch script; adapt resource names to the cluster}
\begin{lstlisting}[language=bash,basicstyle=\ttfamily\scriptsize]
#!/bin/bash
#SBATCH --job-name=torch-r
#SBATCH --nodes=1
#SBATCH --ntasks=1
#SBATCH --cpus-per-task=4
#SBATCH --mem=16G
#SBATCH --time=00:15:00
#SBATCH --gres=gpu:1
#SBATCH --output=logs/%x_%j.out

module load R
# The cluster may also require CUDA or a suitable container/module.
Rscript gpu_torch.R
\end{lstlisting}
\end{block}
\begin{itemize}
  \item Request a GPU explicitly; do not run GPU exercises on a login node.
  \item The Slurm GPU option, partition name, modules, drivers, and supported CUDA version are cluster-specific.
  \item For this four-hour course, one shared demonstration GPU is enough; the main hands-on exercises can remain CPU-based.
\end{itemize}
\end{frame}
